\section{Exact-return languages and length phases}\label{sec:phases55}
The existence of cofinite returns is not the only useful synchronization regime. In fact an exact-return language, when nonempty, has a rigid eventual form. Here and in this section a return has positive length.

\begin{lemma}[The return period]\label{lem:returnperiod55}
Suppose
\[
 \mathcal R=\{|w|>0:u_w=1\}\ne\varnothing,
 \qquad p=\gcd\mathcal R.
\]
There is an integer $g_p$ such that every multiple $dn$ with $n\ge g_p$ belongs to $\mathcal R$, and no length outside $p\N$ belongs to $\mathcal R$. In particular, every nonempty exact-return language is eventually one arithmetic progression.
\end{lemma}
\begin{proof}
Concatenation makes $\mathcal R$ additive. A finite subset of its elements already has gcd $p$: start with one length and successively decrease the gcd when necessary; a strictly decreasing chain of its positive divisors is finite. Divide these finitely many lengths by $p$. Their gcd is one, so the residue argument of Proposition~\ref{prop:return54} shows that their nonnegative sums contain every sufficiently large integer. Multiplying by $p$ proves the assertion. The exclusion of other lengths is the definition of the gcd.
\end{proof}

Fix a letter with product $a\in H$, and let $J$ be the closure of the products of words whose lengths are multiples of $p$, including the empty word.
\begin{lemma}[The phase subgroup]\label{lem:phasegroup55}
The group $J$ is an open normal subgroup of $H$. The quotient $H/J$ is cyclic of order dividing $p$, generated by $aJ$, and
\begin{equation}\label{eq:phasereachable55}
 \overline{\{u_w:|w|\equiv r\pmod p\}}=a^rJ
 \qquad(0\le r<p).
\end{equation}
If $H/H^\circ$ is finite, then $J^\circ=H^\circ$.
\end{lemma}
\begin{proof}
The compact closure defining $J$ is a group, by the positive-power argument, and is generated densely by the block alphabet $\mathcal A^p$. Also $a^p\in J$. If $x=u_w$ for a word $w$ of length $p$, then the executable word consisting of $p-1$ copies of the chosen letter, followed by $w$, followed by one more copy, has product $axa^{p-1}$ and length $2p$. Thus $axa^{p-1}\in J$, whence $axa^{-1}\in J$. It follows that $aJa^{-1}\subset J$. Applying this inclusion $p$ times and using $a^p\in J$ shows equality. Every letter $b$ has $ba^{p-1}\in J$, so $ba^{-1}\in J$. Thus all letters lie in $Ja$, normalize $J$, and their products of length $n$ lie in $Ja^n$. The finite union $\bigcup_{r=0}^{p-1}Ja^r$ is a closed group containing the letters, hence is $H$. This proves normality, finite index, and openness. Conversely, words consisting of $r$ copies of $a$ followed by arbitrary $p$-blocks have products dense in $Ja^r$, proving \eqref{eq:phasereachable55}. An open subgroup contains the identity component, and its own identity component is therefore the same one.
\end{proof}
The quotient order may be a proper divisor of $p$; the phase radii repeat whenever the physical cosets repeat. The phase index records executable length, not merely the physical component quotient.

Assume now that $H$ has finitely many components and put $K=H^\circ$. For $r\in\{0,\ldots,p-1\}$ define
\begin{equation}\label{eq:phaseradius55}
 \varepsilon_r=\max_{s,j,\ c\in H/K:\ c\subset a^rJ}
                    \rad_{Y_j}(F_{s,j}(c)).
\end{equation}
\begin{theorem}[The phasewise clocked threshold]\label{thm:phases55}
Under the sole synchronization assumption $\mathcal R\ne\varnothing$, for every $r$ and $\varepsilon\ge0$,
\[
 \sup_{N\equiv r\ (\mathrm{mod}\ p)}W_{N,\varepsilon}<\infty
             \quad\Longleftrightarrow\quad
             \varepsilon\ge\varepsilon_r.
\]
Equality is attained by a finite component-permutation realization on the indicated horizons. If $\varepsilon<\varepsilon_r$, then for each $k$ there is a constant $L_{k,r}$, independent of $N$, such that every error-$\varepsilon$ realization at $N\equiv r\pmod p$ satisfies
\[
                 \#\{1\le t\le N:|S_t|\le k\}\le L_{k,r}.
\]
The intervening widths are unrestricted. Consequently the full family is bounded exactly at and above $\max_r\varepsilon_r=\varepsilon_{\mathrm c}$; below it, divergence is asserted phase by phase, not necessarily along every horizon.
\end{theorem}
\begin{proof}
Every product of length $N\equiv r$ lies in $a^rJ$. Store its component and the seed and decode component centers as before. This proves the upper bound, including equality, with at most $|\mathcal S||H/K|$ labels. Centers on other components may be chosen arbitrarily when only this phase is under consideration.

For the converse fix a component $c\subset a^rJ$ and $s,j$ with radius greater than $\varepsilon$. We must count narrow cuts in every residue class, not only at block boundaries of one phase. For each $\ell\in\{0,\ldots,p-1\}$ choose $b_\ell\in\{0,\ldots,p-1\}$ congruent to $r-\ell$ modulo $p$. By \eqref{eq:phasereachable55} and the openness of components there is a fixed word $v_\ell$, of length $p_\ell\equiv\ell\pmod p$, with product in $a^{-b_\ell}c$. Indeed this component is contained in $a^\ell J$. Consider the block experiment on $J$ with alphabet $\mathcal A^p$ and target
\[
        G_\ell(h)=F_{s,j}(a^{b_\ell}h u_{v_\ell}),\qquad h\in J.
\]
It is a legal compact-convex interface. Its image on $K=J^\circ$ is $F_{s,j}(c)$, since $K$ is normal. Lemma~\ref{lem:returnperiod55} gives cofinite exact returns in block length for this alphabet.

Restrict the original machine to words that begin with $v_\ell$, end with $b_\ell$ copies of the chosen letter, and have $M=(N-p_\ell-b_\ell)/p$ arbitrary $p$-blocks in between. When $M\ge0$, absorb the fixed prefix into the initialization and the fixed suffix into the terminal decoder; stochastic averaging keeps that decoder in $Y_j$. This produces an error-$\varepsilon$ clocked block machine for $G_\ell$, with registers at original cuts $p_\ell+dm$, $0\le m\le M$. Theorem~\ref{thm:radius54} bounds its narrow cuts $1\le m\le M$ by a constant depending only on $k$ and the fixed data. The omitted cuts of residue $\ell$ number at most $p_\ell+b_\ell+1$: this also charges the block cut zero, while all counted cuts remain in $1,\ldots,N$. If $M<0$, the whole horizon is bounded by $p_\ell+b_\ell$. Sum these bounds over the $p$ residue classes. This counts all cuts $1,\ldots,N$, independently of all other register widths.
\end{proof}
The constants in this theorem depend on the fixed alphabet, return period and certificate, interface, error, width $k$, and chosen finite prefixes. They do not depend on the horizon or on the machine. The original constants $B,d,\chi,L_k$ in Section~\ref{sec:local54} use $d$ for a contraction defect. The return period here is denoted $p$ to keep the two parameters distinct.

\begin{proposition}[Different phases can have different thresholds]\label{prop:parity55}
There is an alphabet with return period two for which $W_{N,0}=1$ on every even horizon, but $W_{N,\varepsilon}\to\infty$ along odd horizons for every $0\le\varepsilon<\rho/2$.
\end{proposition}
\begin{proof}
Take $H=(\R/2\pi\Z)\times\Z/2\Z$ and letters $(\theta,1),(-\theta,1)$, with $\theta/2\pi$ irrational. The generated group is dense in $H$. An exact return must have zero signed count and hence even length; the two-letter return gives every positive even length. Here $J=(\R/2\pi\Z)\times\{0\}$. For one binary-mean query put $F(\phi,0)=0$ and $F(\phi,1)=\rho\cos\phi$, where $0<\rho\le1$. The even-phase target is identically zero and needs one label. The odd-phase radius in the norm $|\cdot|/2$ is $\rho/2$. Apply Theorem~\ref{thm:phases55}.
\end{proof}
